\chapter{Att kontrollera ett påstående}
\label{app:verifiera}

Det här materialet gör, som allt undervisningsmaterial, en rad
påståenden.  De flesta kan ni pröva själva genom att köra programmet ---
det är hela poängen med att koden på bilderna är samma kod som körs, och
att körningen längst bak går att upprepa.  Men några påståenden går inte
att pröva på det sättet:
\begin{enumerate}
  \item att en slumpgenerator ger samma följd av tal om den får samma
    startvärde, och att den utan startvärde utgår från klockan
    (\cref{sec:slump});
  \item att listan, tupeln och de övriga behållarna beter sig som
    materialet beskriver;
  \item de didaktiska påståenden som anteckningarna i marginalen gör om
    hur människor lär sig.
\end{enumerate}
Hur vet man om sådant stämmer?  Samma fråga möter er varje gång ni läser
en lärobok, en webbsida eller ett svar från en språkmodell --- och varje
gång ni själva ska skriva en rapport.  Den här bilagan beskriver en
arbetsgång för att kontrollera ett påstående.
\Cref{tab:guess-pastaenden} visar var vart och ett av materialets
påståenden är belagt.

\section{Gör påståendet till en fråga}

Ett påstående går att kontrollera först när det är tydligt vad som skulle
kunna visa att det är \emph{fel}.  Formulera det därför som en fråga med
ett svar: \enquote{tar det längre tid att söka i en lista ju längre listan
är?} kan besvaras med ja eller nej, medan \enquote{listor är viktiga} inte
kan det --- ingen tänkbar källa skulle visa att det är fel, eftersom
\enquote{viktigt} är ett omdöme och inte ett faktum.  Ofta döljer sig
flera påståenden i ett: \enquote{stegvis förfining är en etablerad metod
som härrör från Wirth} --- ett påstående ni mötte i modulen om
algoritmiskt tänkande --- är två frågor, om \emph{ursprung} och om
\emph{etablering}, och de kräver olika slags källor.

\section{Välj rätt sorts källa}

Olika frågor besvaras av olika källor.
\begin{description}
  \item[Uppslagsverk] för vad ett ord \emph{betyder}.
  \item[Primärkällan] för var en idé \emph{kommer ifrån}, och för hur ett
    system \emph{beter sig}.  Frågan om vad en av Pythons behållare gör
    besvaras av Pythons egen dokumentation, inte av en lärobok som
    återger den.
  \item[Översikter] (\foreignlanguage{english}{reviews}) för om något är
    \emph{etablerat}, eller för vad forskningen \emph{sammantaget} säger.
  \item[Enskilda studier] för ett specifikt resultat --- med förbehållet
    att en enskild studie bara visar att något observerats en gång.
\end{description}
Var man hittar dem: universitetsbibliotekets söktjänst, och databaser som
Scopus, Web of Science, IEEE Xplore, ACM Digital Library, DBLP (datalogi)
och OpenAlex (öppen, och trots namnet inte detsamma som öppen tillgång).
Google Scholar täcker brett men rankar ogenomskinligt.

\section{Sök så att du kan ha fel}

Den som söker på namn hon redan känner till hittar det hon väntar sig.
Sök därför \emph{ämnesdrivet} --- på begrepp, inte författare.  Några
praktiska regler: håll frågorna korta, ställ samma frågor i flera
databaser, sök begreppets alla namn, och sök också efter
\emph{invändningar} genom att lägga till ord som \enquote{critique},
\enquote{limitations} eller \enquote{failure}.  Ett påstående som
överlevt en motsökning är värt mer än ett som aldrig prövats, och en
invändning som hittas blir ett \emph{förbehåll} i svaret, inte ett
nederlag.

\section{Läs i fulltext, och skriv ner hur}

En träff i en databas är inte ett belägg, och sammanfattningen
(\foreignlanguage{english}{abstract}) räcker inte heller: den säger vad
författarna vill ha sagt, inte vad de visat.  Läs hela texten och leta
upp \emph{det ställe} som styrker påståendet --- ett ordagrant citat med
sidnummer.  Räkna aldrig en källa ni inte läst.

Skriv sedan ner hur ni gick till väga, så att någon annan kan följa
spåret.  I det här materialet står den anteckningen som en kommentar
ovanför varje post i litteraturlistans källfil, med fasta fält:
\begin{description}
  \item[\texttt{CLAIM}] vilket påstående källan ska styrka;
  \item[\texttt{FOUND-VIA}] hur den hittades;
  \item[\texttt{PICKED}] varför just den, framför andra träffar;
  \item[\texttt{QUOTE}] det ordagranna citatet, med sida eller avsnitt;
  \item[\texttt{VERIFIED}] att fulltexten lästs, och var;
  \item[\texttt{COUNTER}] vad motsökningen gav;
  \item[\texttt{DATE}] när.
\end{description}
Det tar ett par minuter per källa, och det är det enda som gör det
möjligt att om ett år svara på frågan \enquote{hur vet du det?}

\section{Dra slutsatsen --- med förbehåll}
\label{sec:sok-slutsats}

Svaret skrivs ut tillsammans med det som talar emot.  Ett påstående om
hur snabb en operation är gäller till exempel en viss implementation och
ett genomsnittsfall, inte alla tänkbara fall --- och det förbehållet hör
till svaret.  Förbehållen är inte svagheter i svaret.  De \emph{är}
svaret, och de avgör hur påståendet får användas i materialet.

\section{En anmärkning om språkmodeller}

En språkmodell kan användas för att \emph{hitta} och \emph{sortera}
källor.  Lägg märke till arbetsfördelningen: varje källa som citeras ska
läsas och väljas av en människa.  Använd gärna språkmodeller för att
hitta och sortera; låt dem aldrig vara det sista ledet.  De kan göra
fel, men det är ni själva som måste stå till svars för innehållet.

\section{Materialets påståenden och var de är belagda}
\label{sec:guess-pastaenden}

Inget av den här modulens påståenden krävde en egen litteratursökning.
Frågor om vad Python gör besvaras av Pythons egen dokumentation, som är
primärkällan: hur behållarna beter
sig\autocite{PythonDataStructures} och hur slumpgeneratorn
fungerar\autocite{PythonRandom}.  Den sista säger både det materialet
använder --- \enquote{The generator's \mintinline{python}{random()}
method will continue to produce the same sequence when the compatible
seeder is given the same seed} --- och det förbehåll som hör till:
generatorn är \enquote{completely deterministic} och därför
\enquote{completely unsuitable for cryptographic purposes}.  Materialet
återger båda delarna, eftersom ett startvärde som gör en körning
upprepbar i ett läromedel är samma egenskap som gör generatorn
olämplig när slumpen ska vara oförutsägbar.

Principerna om ett ansvar per funktion
(\foreignlanguage{english}{SRP}) och om att hålla lösningen enkel
(\foreignlanguage{english}{KISS}) namnges här men beläggs i kursens
modul om funktioner.  Att kontrast måste upplevas samtidigt för att en
aspekt ska kunna urskiljas är hämtat från Martons \emph{Necessary
Conditions of Learning}\autocite[45--47]{Marton2015}.

\begin{table}[htbp]
  \centering
  \caption{Modulens påståenden, som frågor, och var svaret finns.}
  \label{tab:guess-pastaenden}
  \begin{tabular}{@{}p{0.46\linewidth}p{0.44\linewidth}@{}}
    \toprule
    Fråga & Var det är belagt \\
    \midrule
    Ger samma startvärde samma följd av slumptal, och vad händer utan
    startvärde?
      & Primärkälla: Pythons dokumentation för modulen
        \mintinline{python}{random}, läst i fulltext 2026-09-03; körningen
        i \cref{sec:slump} går dessutom att upprepa. \\
    Beter sig behållarna som materialet beskriver?
      & Primärkälla: Pythons egen dokumentation, läst i fulltext
        2026-09-03. \\
    Ska en funktion ha ett ansvar, och ska lösningen hållas enkel?
      & Belagt i föreläsningen \emph{Funktioner}; namnen återanvänds här. \\
    Måste kontrast upplevas samtidigt för att en aspekt ska kunna
    urskiljas?
      & Primärkälla: Marton, \emph{Necessary Conditions of Learning},
        s.~45--47. \\
    \bottomrule
  \end{tabular}
\end{table}
